\section{Conclusion}
\label{sec:conclusion}

The central finding of this work is that automated element identification in CF-LIBS at low resolving power (RP~$= 300$--$1100$) is fundamentally limited by the spectral information content available to any single identification method, and that combining complementary spectral evidence---global Boltzmann emissivity fitting, individual line-position matching, multiplet spacing recognition---through systematic optimization yields identification performance ($F_{\beta=0.5} = 0.914$) that far exceeds what any individual method can achieve.
GPU acceleration is the enabling technology that makes both the systematic benchmarking and the exhaustive multi-method optimization practical.
We have presented three co-equal contributions to the advancement of Calibration-Free LIBS: a GPU-accelerated computational pipeline, a systematic benchmark of element identification methods at low resolving power, and a GPU-optimized multi-method consensus identification system that dramatically outperforms any individual approach.
Built on the JAX framework and benchmarked on an NVIDIA Tesla V100S GPU, the first contribution comprises five GPU-optimized core algorithms:

\begin{enumerate}
  \item The \textbf{Voigt profile} kernel, based on the branch-free \citet{Weideman1994} $N\!=\!36$ rational approximation to the Faddeeva function, achieves a $76.4\times$ throughput increase at $10^5$ wavelength grid points with a maximum relative error of $5.1 \times 10^{-14}$ against the SciPy reference.
  % Source: benchmark_summary.json/headline_numbers/voigt_max_speedup = 76.38
  % Source: voigt_results.json/results/max_relative_error ~ 5.07e-14

  \item The \textbf{vectorized Boltzmann fitting} kernel uses closed-form weighted least squares with a pad-and-mask strategy, achieving $8.8\times$ speedup at 20~elements with near-constant GPU execution time regardless of element count.
  % Source: benchmark_summary.json/headline_numbers/boltzmann_max_speedup = 8.83

  \item The \textbf{Anderson-accelerated Saha--Boltzmann solver} reduces the average iteration count by $1.6\times$ at tolerance $10^{-6}$, with depth $m = 1$ optimal for typical LIBS conditions and deeper history ($m = 5$) providing $2.9\times$ reduction at tighter tolerances.
  % Source: benchmark_summary.json/headline_numbers/anderson_avg_iteration_reduction = 1.62

  \item The \textbf{softmax closure} reparameterization provides an analytically differentiable mapping from unconstrained parameters to the compositional simplex, enabling gradient-based joint optimization of temperature, electron density, and elemental concentrations.

  \item The \textbf{batch forward model} achieves a peak throughput of $10{,}708$ spectra per second with bit-identical results to sequential computation, and an end-to-end pipeline speedup of $13.6\times$ at batch size 1000.
  % Source: benchmark_summary.json/headline_numbers/batch_forward_max_throughput_spectra_per_sec = 10708.1
  % Source: benchmark_summary.json/headline_numbers/e2e_max_speedup = 13.64

\end{enumerate}

\noindent The second contribution is a \textbf{systematic benchmark of five element identification pathways} on 74~Aalto spectra at RP~$= 300$--$1100$, establishing quantitative performance baselines.
Per-algorithm standalone evaluation reveals complementary strengths---NNLS achieves the highest recall ($R = 0.922$), the hybrid NNLS+ALIAS achieves the best single-method precision ($P = 0.578$), and no individual algorithm dominates---with per-element analysis identifying Mn and Na as fundamental precision limiters at RP~$< 1000$.

The third contribution is \textbf{GPU-accelerated multi-method consensus identification optimization} over a 101-parameter space spanning all five identification methods, achieving $F_{\beta=0.5} = 0.914$---a 121\% improvement over the hand-tuned baseline ($F_{\beta} = 0.413$)---by learning per-element identification confidence thresholds, method selection weights, and evidence combination parameters.
The optimization evaluates $2 \times 10^6$ configurations in 58~seconds on a single V100S GPU, demonstrating that GPU acceleration enables exhaustive systematic exploration of multi-method identification configurations.

These findings demonstrate that (1)~GPU acceleration is viable for the full CF-LIBS pipeline with no loss of numerical accuracy; (2)~no single identification method dominates at low RP, because each exploits different spectral information; and (3)~GPU-accelerated optimization enables exhaustive exploration of a 101-dimensional identification configuration space that would be infeasible manually.
Validation on 74~real LIBS spectra and 6~ChemCam calibration targets confirms machine-precision GPU--CPU parity across all kernels.
The GPU becomes faster than CPU at batch sizes $\geq 10$, establishing a clear operating regime for real-time deployment.
% Source: benchmark_summary.json/headline_numbers/e2e_crossover_batch_size = 10

Several directions merit future investigation.
On the computational side, multi-GPU scaling via JAX's \texttt{pmap} would enable manifold generation over parameter grids exceeding $10^6$ plasma conditions with linear throughput scaling, while a mixed-precision approach using float32 for the profile matrix and float64 for the Weideman coefficients could approximately double throughput for applications where $\sim 0.1\%$ accuracy is acceptable (e.g., qualitative elemental screening or rapid survey campaigns).
Accelerating the nearest-neighbor search in the spectral manifold with GPU FAISS would complete the GPU-accelerated chain from raw spectrum to elemental composition.
Extending the physical model to include non-LTE corrections via collisional-radiative models would broaden the range of analyzable plasma conditions beyond the partial thermodynamic equilibrium assumed here, enabling analysis of early-gate plasmas and high-ionization-potential elements.
On the deployment side, porting the pipeline to embedded GPU platforms (NVIDIA Jetson, Apple Neural Engine) would enable field-deployable real-time LIBS analysis for industrial process monitoring and geological survey applications.
For the element identification system, cross-instrument evaluation of the optimized 101-parameter configuration on spectra from diverse spectrometer types (different RP, wavelength ranges, detector technologies) would establish generalization bounds and determine whether the learned per-element thresholds transfer across instruments or require per-instrument re-optimization.
Benchmarking at higher resolving powers (RP~$\approx 5000$, e.g., Mechelle~5000 echelle spectrometer) would quantify the RP-dependent identification specificity curve and establish the upper performance bound for peak-matching approaches.

The source code, benchmark scripts, and validation data are publicly available to facilitate reproducibility and adoption by the LIBS community.
